\section{Discussion}
\label{sec:discussion}

We discuss the implications of our results for silicon-sampling methodology, for the broader use of LLMs as representation generators, and for survey practice, before turning to limitations.

\subsection{Why Representation-Based Silicon Sampling Works}
The central empirical finding we anticipate---that a representation-based factorization outperforms direct generation, especially under transfer---admits a structural explanation. Direct generation forces the LLM to commit to a single answer token, which is a lossy readout of a high-dimensional latent reaction distribution; the token conflates respondents whose distributions differ in their lower-ranked options, inherits the model's response bias, and couples the model's world knowledge to the specific instrument phrasing. By reading out a vector instead, we preserve the geometry of the reaction distribution, and the option-aware decoder can then anchor this geometry to the concrete option set of the target instrument and be calibrated against its empirical choice distribution. This decoupling---between the representation of the respondent's reaction and the representation of the instrument's options---is the survey-domain analogue of the decoupling between location semantics and temporal model that \citet{he2025geolocation} exploit, and it is what enables the cross-instrument and cross-wave transfer we expect to observe. The use of a generative embedder~\citep{behnamghader2026llmvecgen} is essential here: a discriminative encoder optimized for retrieval would discard the generative disposition semantics that make the reaction vector a faithful surrogate for the answer distribution.

\subsection{The Boundary Between Query Conditioning and Reuse}
Our second contribution is the explicit characterization of the boundary between the query-conditioned ResponseVec and the reusable RespondentVec. The decomposition of Section~\ref{sec:tradeoff} predicts that the per-item accuracy gap $\Delta(q)$ is positive for high-discriminability items within an instrument but shrinks or reverses under transfer, and our anticipated Pareto analysis (Section~\ref{sec:results-q3}) operationalizes this by showing that a mixed allocation policy dominates uniform strategies. This finding has a practical implication: survey practitioners imputing large panels should not choose uniformly between the two regimes but should allocate ResponseVec to items where the item signal exceeds the overfitting penalty, estimated on a small calibration set. The shift of the Pareto frontier between the unseen-item and OOD-intersection regimes further implies that the optimal allocation is itself regime-dependent, with OOD-intersection imputation favoring the cheaper RespondentVec more heavily. This is, to our knowledge, the first actionable cost--accuracy guidance for individual-level silicon sampling.

\subsection{Implications for Survey Research}
Beyond the methodological advance, our results bear on the broader debate over LLMs in survey research. \citet{westwood2025potential} frames LLMs as a potential existential threat to online surveys; our work suggests a more constructive role, in which LLMs augment rather than replace human respondents by imputing missing items in panel data, projecting opinions onto unobserved demographic intersections, and enabling counterfactual reasoning about unasked questions---all while preserving calibrated uncertainty. The calibration advantage of the representation-based approach is particularly important here: a silicon-sampled imputation that is well-calibrated can be used for downstream statistical inference (e.g., estimating opinion gaps with honest confidence intervals), whereas an uncalibrated direct-generation imputation would produce misleadingly precise estimates. The OOD-intersection transfer regime (R3) is especially relevant for this use case, and our anticipated finding that RespondentVec is the most OOD-robust method suggests that stable-disposition encodings are the right tool for imputation onto populations underrepresented in training.

\subsection{Limitations}
Our study has several limitations. First, our evaluation uses the ISSP-based SocioBench benchmark across four domains and roughly 30 countries; while the multi-country structure supports cross-demographic transfer, generalization to non-ISSP instruments or to genuine longitudinal panels requires further validation, and the well-documented WEIRD bias of LLMs~\citep{bail2024generative} may constrain transfer to other populations. Second, the generative embedder LLM2Vec-Gen is itself trained on internet text that may contain survey-like content, raising a subtle contamination concern; we mitigate this by using off-the-shelf encoders without survey-specific continued training in our main results, but a thorough contamination audit is left for future work. Third, our cost analysis measures GPU-hours and latency but not dollar cost, which depends on deployment infrastructure; the relative ordering of methods should be robust, but absolute budgets will differ. Fourth, our ablations isolate design choices one at a time, leaving interactions (e.g., between encoder choice and decoder design) for future study. Fifth, we do not implement a separate generative-agent (GEMS) baseline; as noted in Section~\ref{sec:baselines}, its memory-stream architecture collapses under our canonical prompt into the strong-prompting path already covered by B1 with history, and a faithful re-implementation would constitute a separate study. Finally, our anticipated results are framed as expectations consistent with the protocol; the empirical confirmation of predictions P1--P3 is the burden of the experimental runs, and any deviation would itself be informative about the limits of the representation-based factorization.

\subsection{Ethical Considerations}
Individual-level silicon sampling raises dual-use concerns: a method that accurately predicts how a specific person will answer survey questions could be misused for microtargeting or for de-anonymizing panel data. We mitigate this in two ways. First, our method requires access to the respondent's own prior answer history, which is not available for arbitrary individuals; it is a panel-imputation tool rather than a population-inference tool. Second, we frame our work throughout as augmentation of existing survey data with calibrated uncertainty, not as replacement of human respondents, and we support the call of \citet{westwood2025potential} for survey-methodology safeguards against LLM-based contamination. We will release code and configurations but not any individual-level survey data, which remains governed by the terms of the ISSP/SocioBench data-use agreements.
